\documentclass[11pt, letterpaper]{article}

\usepackage[utf8]{inputenc}
\usepackage[T1]{fontenc}
\usepackage{mathpazo} % Palatino font for a literary, magazine feel
\usepackage{geometry}
\geometry{left=1.25in, right=1.25in, top=1.25in, bottom=1.25in}
\usepackage{setspace}
\onehalfspacing
\usepackage{titlesec}
\usepackage{hyperref}

\titleformat{\section}{\large\scshape\centering}{}{0em}{}

\title{\textbf{The Illusion of Sanctuary:\\Why Canada Must Stop Hiding Its Failures Behind MAiD Restrictions}}
\author{Justin Bogner}
\date{}

\begin{document}

\maketitle

\vspace{-2em}
\begin{center}
\textit{Prepared for Submission to The Walrus}
\end{center}
\vspace{2em}

In February 2022, a fifty-one-year-old Ontario woman known publicly only as Sophia chose medical assistance in dying. She lived with multiple chemical sensitivities so acute that cigarette smoke, cleaning products, and everyday building materials rendered most housing toxic to her. For more than two years she and her advocates begged provincial and federal officials for a single safe, affordable apartment. None was provided. On the day her MAiD procedure was approved, her suffering was described as intolerable---not solely from her medical condition, but from the cumulative weight of systemic abandonment. Sophia became one of the earliest and most visible symbols of a quiet crisis: Canadians turning to the state's ultimate medical service because the state had failed to furnish the most basic conditions of a livable life.

Her story is no outlier. In 2024, Health Canada recorded 16,499 MAiD deaths---5.1 percent of all deaths in the country. Track 2 cases (those where natural death is not reasonably foreseeable) continue to rise, now numbering in the hundreds annually. Disability advocates and independent reports document a recurring pattern: requests driven less by untreatable physical decline than by the grinding realities of legislated poverty, inaccessible housing, and absent community care. Provincial disability supports often max out between \$800 and \$1,800 per month---well below the poverty line in every major city---while waitlists for accessible or chemically safe housing stretch years or decades. When these structural failures intersect with serious illness or disability, the choice to die can begin to look less like a free exercise of autonomy and more like the only exit left on the table.

An unlikely coalition has formed in response. Conservative politicians wary of expanding assisted dying and progressive disability-rights organizations alarmed by the devaluation of disabled lives now speak in near-unison: in a country that cannot house its citizens or pay them enough to survive with dignity, offering a medical exit to the marginalized is not compassion. It is coercion by neglect. A person who cannot secure stable housing or adequate income is not exercising sovereign choice when they request MAiD; they are navigating a landscape deliberately engineered to make endurance unbearable. To authorize their death under these conditions, the argument runs, is for the state to ratify the consequences of its own austerity.

Faced with this indictment, Ottawa has chosen delay. The planned expansion of MAiD to individuals whose sole underlying condition is a mental illness---originally slated for 2024, then pushed to 2026---has now been deferred until March 17, 2027. Parliament has struck a special joint committee to study ``readiness,'' while private members' bills circulate to block the change permanently. The government frames these pauses as prudent safeguards. In reality, they constitute a sophisticated form of institutional entrapment.

The logic is as elegant as it is cruel. \textit{We acknowledge that our social safety net has failed you. Because that failure has made your life intolerable, we cannot trust your request to die. Therefore, to protect you from the results of our neglect, we will compel you to keep living in the very conditions that drove you to seek release.} This is not protection. It is paternalism weaponized against the vulnerable. Denying MAiD does not magically produce subsidized apartments, living wages, or timely home care. It simply removes the one option the state has made available, leaving individuals trapped in the precise state of despair that prompted their request. The policy converts the government's own policy failures into a perpetual life sentence, all while shielding politicians from the political cost of visible deaths.

The disability-rights movement is correct that systemic deprivation coerces. The United Nations has warned that expanding MAiD without first guaranteeing adequate supports risks violating the Convention on the Rights of Persons with Disabilities. Yet the answer to state neglect cannot be the further restriction of fundamental liberties. We do not remedy a housing crisis by barricading the exits and holding suffering citizens hostage to our moral discomfort. We do not vindicate the right to life by stripping away the right to refuse a life the state has rendered unlivable.

Canada needs a \textbf{Sovereignty Model} that treats decisional capacity as the sole threshold for eligibility while imposing rigorous structural accountability on the state itself.

Under this framework, MAiD eligibility would rest exclusively on whether an individual possesses the cognitive and emotional capacity to understand their diagnosis, prognosis, available treatments, and the irreversible nature of the choice. Capacity assessments would be performed by independent, multidisciplinary teams (physicians, psychiatrists, ethicists) rather than the patient's usual care providers, with mandatory cooling-off periods and documented exploration of all reasonable alternatives. For psychiatric conditions, additional safeguards---such as documentation that evidence-based treatments have been exhausted or refused---would apply, consistent with practices in the Netherlands and Belgium, where psychiatric MAiD has operated for years under strict protocols without evidence of widespread abuse.

Crucially, while capacity remains the ultimate threshold, the process requires structural transparency when social determinants are the dominant driver. Before any provision of MAiD, an independent resource-screening body would be required to certify that the state has offered concrete, time-bound supports: a housing placement, an increase in disability income to a livable level, intensive community care, or whatever specific intervention the individual identifies as necessary to make continued life tolerable. If those supports do not exist---if the provincial waitlist for accessible housing is measured in years, or if disability stipends remain anchored below the poverty line---the state would be legally obligated to document that failure in a public registry.

Health Canada would then publish an annual, disaggregated ledger: how many MAiD recipients cited housing insecurity, inadequate income, or social isolation as primary or contributing factors; which provinces generated the highest proportions; and what concrete supports were (or were not) available at the time of request. Imagine the political pressure if Parliament were forced, each spring, to confront a documented body count of citizens who chose death specifically because their governments failed to provide shelter or a living wage. The Sovereignty Model transforms private tragedy into public indictment. It refuses to let the state hide its failures behind the closed doors of clinical encounters while simultaneously denying citizens the dignity of self-determination.

This approach does not dismiss the legitimate fears of disability advocates. It incorporates them. By making social coercion visible and quantifiable, it creates the precise political incentive structure needed to force provincial and federal governments to fund the social determinants of health---housing, income security, home care---at the scale the crisis demands. Paternalistic gatekeeping, by contrast, merely sweeps the victims back under the rug and allows politicians to claim moral high ground while doing nothing.

The philosophical stakes are straightforward. In \textit{Carter v. Canada} (2015), the Supreme Court recognized that the prohibition on assisted dying violated section 7 of the Charter---the right to life, liberty, and security of the person---precisely because it forced individuals to endure intolerable suffering against their will. That liberty interest does not evaporate when the suffering is compounded by poverty or policy failure. Self-ownership is not a privilege granted only to the comfortably housed and well-supported. To condition it on the state's prior perfection of the welfare system is to hold citizens hostage to a utopia that may never arrive.

Canada's current path---endless study, repeated delays, and the comforting fiction that restriction equals compassion---serves neither autonomy nor justice. It preserves the illusion of sanctuary while the real sanctuary, the one built from adequate housing, income, and care, remains unbuilt. The Sovereignty Model offers a different bargain: absolute respect for the individual's ultimate jurisdiction over their own body and mind, paired with an unflinching obligation on the state to answer, in public and in numbers, for every life it has made unbearable.

The choice before Parliament in 2027 is not between compassion and cruelty. It is between paternalistic entrapment and genuine respect for persons as sovereign agents. Sophia's death, and the thousands like it that statistics now quietly record, demand that we choose the latter. Canada must stop using MAiD restrictions as a moral smokescreen for its crumbling social contract. The vulnerable deserve better than forced endurance. They deserve the dignity of choice---and a state finally compelled to make that choice a last resort rather than the only one left.

\end{document}
